\section{Signorini's Problem}


\subsection{Physical and Geometric Setup of the Problem}

We consider an elastic body occupying a reference configuration $\Omega\subset\R^d,\,d\in\{2,3\}$.
Under the action of external forces and possible contact with a rigid foundation, the body deforms.
In the {\em small--strain} framework, the deformation is described by the {\em displacement field}
$u:\Omega\to\R^d$, and a material point $x\in\Omega$ is mapped to its deformed position
\begin{equation*}
	\varphi(x) := x' := x + u(x),\quad x\in\overline{\Omega}.
\end{equation*}
The aim of Signorini's problem is to determine an {\em equilibrium} displacement $u$ under
standard boundary conditions and {\em unilateral} contact constraints on the part of the boundary
that may touch the rigid foundation. Throughout the whole report, $\Omega\subset \R^d,\, d\in\{2,3\}$ denotes an open, bounded and connected {\em Lipschitz} domain. The later means, that the boundary $\Gamma:=\partial\Omega$ can locally be described as the graph of a Lipschitz continuous function. This ensures for example, that the outward normal vector $\mathrm{n}(x)$ is well--defined almost everywhere. The boundary is decomposed into three pairwise disjoint parts
\begin{equation*}
	\Gamma = \overline{\Gamma_D}\,\mathbin{\dot\cup}\,\overline{\Gamma_N}\,\mathbin{\dot\cup}\,\overline{\Gamma_C}
\end{equation*}
which are assumed to be open in $\Gamma$ and where each part is subject to different boundary conditions to model different physical interactions (See Figure \ref{fig:setup}). On $\Gamma_D$, we apply {\em homogeneous Dirichlet} boundary conditions. This forces the body to be {\em clamped} and this part, i.e.\ the displacement should not act on $\Gamma_D$,
\begin{equation*}
	\forall x\in\Gamma_D:\;x + u(x) = x \quad\Longrightarrow \quad\trace(u)|_{\Gamma_D} = 0.
\end{equation*}
The part $\Gamma_N$ may be subject to a {\em Neumann} boundary condition, which means we impose a known {\em surface force density}. $\Gamma_C$ denotes the {\em potential contact} boundary. We assume that $\Gamma_C$ has positive $(d-1)$--Lebesgue measure. The rigid foundation is  positioned so that any actual contact, both initially and in the deformed equilibrium configuration, may only occur on $\Gamma_C$. While $\Gamma_C$ is prescribed a priori, the {\em active contact zone} $\Gamma_{\mathrm{act}}(u)\subset\Gamma_C$ is unknown and depends on the solution $u$. Later on, the boundary conditions imposed on $\Gamma_C$ enforce {\em unilateral} contact through {\em nonpenetration} and {\em complementarity}.

\begin{figure}[!h]
	\centering
	\begin{tikzpicture}[scale=1.05, >=Latex, line cap=round, line join=round]
	\definecolor{cD}{RGB}{220,50,47}
	\definecolor{cN}{RGB}{38,139,210}
	\definecolor{cC}{RGB}{42,161,152}
	\coordinate (A) at (0.2,1.35);
	\coordinate (B) at (6.0,1.35);
	\coordinate (D) at (5.6,3.85);
	\coordinate (E) at (1.0,3.65);
	\path[fill=gray!10]
	(A) .. controls (2.0,0.95) and (4.2,0.95) .. (B)
	.. controls (6.55,2.05) and (6.25,3.15) .. (D)
	.. controls (4.55,4.55) and (2.15,4.35) .. (E)
	.. controls (0.25,3.05) and (0.05,2.05) .. cycle;

	\draw[black!55, line width=0.9pt]
	(A) .. controls (2.0,0.95) and (4.2,0.95) .. (B)
	.. controls (6.55,2.05) and (6.25,3.15) .. (D)
	.. controls (4.55,4.55) and (2.15,4.35) .. (E)
	.. controls (0.25,3.05) and (0.05,2.05) .. cycle;
	\draw[cC, very thick]
	(A) .. controls (2.0,0.95) and (4.2,0.95) .. (B);
	\draw[cN, very thick]
	(B) .. controls (6.55,2.05) and (6.25,3.15) .. (D)
	.. controls (4.55,4.55) and (2.15,4.35) .. (E);
	\draw[cD, very thick]
	(E) .. controls (0.25,3.05) and (0.05,2.05) .. (A);
	\node at (3.1,2.45) {$\Omega$};
	\node[cC] at (3.1,1.35) {$\Gamma_C$};
	\node[cN] at (5.55,4.30) {$\Gamma_N$};
	\node[cD] at (-0.35,2.55) {$\Gamma_D$};
	\def\yF{1.024}
	\draw[black, line width=0.9pt] (-0.6,\yF) -- (6.6,\yF);
	\draw[pattern=north east lines] (-0.6,\yF) rectangle (6.6,\yF-0.80);
	\node[below] at (3.0,\yF-0.80) {rigid foundation};
\end{tikzpicture}
	\caption{2D example of Signorini's problem before (left) and after displacement (right).}
	\label{fig:setup}
\end{figure}

\subsection{Weak Formulation}
\subsubsection{Functional Setting and Admissible Displacements}
We start by introducing the energy space and the admissible set of displacements.
Homogeneous Dirichlet boundary conditions on $\Gamma_D$ are incorporated by
\begin{equation*}
	V:=\left\{ v\in H^1(\Omega;\R^d)\;\middle|\; \trace(v)=0 \text{ on }\Gamma_D\right\}.
\end{equation*}
Let $\gamma=\trace:H^1(\Omega;\R^d)\to H^{1/2}(\Gamma;\R^d)$ denote the trace operator.
For $v\in V$ we define the normal component of the boundary displacement on $\Gamma_C$ by
\begin{equation*}
	v_{\mathrm n} := (\gamma v)\cdot \mathrm n \quad \text{on }\Gamma_C,
\end{equation*}
where $\mathrm n$ is the outward unit normal on $\Gamma$. The unilateral nature of contact is expressed by the {\em nonpenetration} constraint. The rigid foundation is located on the exterior side of $\Gamma_C$ in direction of $\mathrm n$. Hence the body is not allowed to move across $\Gamma_C$ into the obstacle. In the small--strain setting this yields
\begin{equation*}
	u_{\mathrm n}\le 0 \quad \text{a.e.\ on }\Gamma_C.
\end{equation*}
This motivates the nonempty, closed and convex cone (\cite[Prop. 3.1]{book}) of admissible displacements
\begin{equation*}
	K := \left\{ v\in V\;\middle|\; v_{\mathrm n}\le 0 \text{ a.e.\ on }\Gamma_C\right\}\subset V.
\end{equation*}
On $V$, we define the symmetric, bilinear and continuous form
\begin{equation*}
	a(u,v):= \int_\Omega \sigma(u):\varepsilon(v)\;\dx,\quad u,v\in V
\end{equation*}
representing the {\em virtual work of internal forces} and the linear bounded functional
\begin{equation*}
	\ell(v):= \int_\Omega f\cdot v\;\dx  + \int_{\Gamma_N} F\cdot v\;\ds,\quad v\in V,\,f\in L^2(\Omega;\R^d),\, F\in L^2(\Gamma_N;\R^d).
\end{equation*}
with {\em volume force density} $f\in L^2(\Omega;\R^d)$ on the body and {\em surface loads} $F\in L^2(\Gamma_N;\R^d)$ on the Neumann boundary.

\subsubsection{Energy Minimization}

The total {\em mechanical energy} associated with the elastic body is given by the functional
\begin{equation*}
	\cJ(v):= \frac{1}{2}a(v,v)-L(v),\quad v\in V.
\end{equation*}



\subsubsection{Variational inequality as an optimality condition}
\begin{proposition}
	Let $V$ be a real Hilbert space, $\emptyset\neq K\subset V$ closed and convex and let $\ell\in L^*$ be a linear and bounded functional. Furthermore, let $a:V\times V\to\R$ be a symmetric bilinear form that is {\em coercive} on $V$, i.e. there exists $\alpha >0$ such that
	\begin{equation*}
		a(v,v)\ge \alpha\|v\|_V^2\quad \forall v\in V.
	\end{equation*}
	Then $u\in K$ fulfills
	\begin{equation}\label{eq:varineq}
		a(u,v-u)\ge \ell(v-u)\quad \forall v\in K
	\end{equation}
	if and only if $u\in K$ is a minimizer of the corresponding energy functional, i.e.
	\begin{equation}\label{eq:energymin}
		u\in \argmin_{v\in K}\cJ(v) := \argmin_{v\in K} \frac{1}{2}a(v,v)-\ell(v).
	\end{equation}
\end{proposition}
\begin{proof}
	Let $u\in K$ fulfill \eqref{eq:energymin} and let $v\in K$ be arbitrary. Since $K$ is convex,
	\begin{equation*}
		u_t:= u+t(v-u)\in K\quad\forall t\in[0,1].
	\end{equation*}
	Set $J(t):= \cJ(u_t)$. Then $J$ has a minimum at $t = 0$ since $u_0 = u$ minimizes $\cJ$. Thus we have $J^\prime(0^+)\ge 0$ and therefore
	\begin{equation*}
		J^\prime(t) =  a(u+t(v-u),v-u)-\ell(v-u)\quad\Longrightarrow\quad J^\prime(0) = a(u,v-u)-L(v-u)\ge 0.
	\end{equation*}
	Now assume that $u\in K$ fulfills \eqref{eq:varineq} and set $w:= v-u$ for some arbitrary $v\in K$. Then
	\begin{equation*}
		\cJ(v)-\cJ(u) = a(u,w)+\frac{1}{2}a(w,w)-L(w)\quad\text{and}\quad a(u,w)-L(w)\ge 0
	\end{equation*}
	which directly implies
	\begin{equation*}
		\cJ(v)-\cJ(u)\ge \frac{1}{2}a(w,w)\ge 0.
	\end{equation*}
	Therefore $\cJ(u)\le \cJ(v)$ for every $v\in K$ which makes $u\in K$ a minimizer of $\cJ$.
\end{proof}
\subsubsection{Existence and Uniqueness}
\subsection{Strong Formulation}
\subsection{Mixed Formulation}
\subsubsection{Mixed Weak Foram as Sattlepoint}
\subsubsection{Well--posedness and inf--sup condition}